%% ma: L12-B14-0074
%% lop: 12
%% chuong: 5
%% bai: 14
%% dang: 12.14.21
%% loai: TN4
%% muc: TH
%% dap_an: "A"
%% nguon: "To_n_12_PT_M_t_Ph_ng.pdf (Đề kiểm tra đánh giá kiến thức, Chuyên đề Phương trình mặt phẳng, Nguyễn Anh Khôi - Lê Minh Khôi), Câu 29"
%% kien_thuc: "Bài 14 (phương trình mặt phẳng theo đoạn chắn), Bài 7 (hình chiếu của điểm trên trục tọa độ)"
%% trang_thai: chinh-thuc
%% ghi_chu: "Tài liệu gốc không kèm đáp án; đáp án do Claude giải và kiểm bằng sympy. Dạng đã duyệt 07/11/2026."
\begin{ex}
Trong không gian $Oxyz$, cho điểm $M(1;2;3)$. Gọi $A$, $B$, $C$ lần lượt là hình chiếu vuông góc của điểm $M$ lên các trục $Ox$, $Oy$, $Oz$. Viết phương trình mặt phẳng $(ABC)$.
\choice{\True $\dfrac{x}{1}+\dfrac{y}{2}+\dfrac{z}{3}=1$}{$\dfrac{x}{1}-\dfrac{y}{2}+\dfrac{z}{3}=1$}{$\dfrac{x}{1}+\dfrac{y}{2}+\dfrac{z}{3}=0$}{$-\dfrac{x}{1}+\dfrac{y}{2}+\dfrac{z}{3}=1$}
\loigiai{$A(1;0;0)$, $B(0;2;0)$, $C(0;0;3)$ nên $(ABC)\colon \dfrac{x}{1}+\dfrac{y}{2}+\dfrac{z}{3}=1$. Chọn A.}
\end{ex}
